\section{Method}

\subsection{Problem Formulation}

Cancer does not progress in isolation. A tumor cell's fate depends critically on signals from its microenvironment—the fibroblasts remodeling the matrix, the macrophages that may nurture or attack, the T cells recognizing antigens or succumbing to exhaustion. This biological reality motivates our central hypothesis: cross-sectional progression signal becomes more identifiable when conditioned on receiver-centered local niche context.

Let $\mathcal{D} = \{(x_i, s_i, d_i, \mathcal{N}_i)\}_{i=1}^{N}$ denote a dataset of $N$ cells, where $x_i \in \mathbb{R}^{G}$ is the log-normalized expression over $G$ genes, $s_i \in \{\text{Normal}, \text{AAH}, \text{AIS}, \text{MIA}, \text{LUAD}\}$ is the histopathological stage, $d_i$ identifies the donor, and $\mathcal{N}_i = \{(x_j, r_{ij}, \tau_j)\}$ represents the spatial neighborhood with distances $r_{ij}$ and cell types $\tau_j$.

Our goal is to learn a representation $f_\theta$ mapping a receiver cell and its neighborhood to an embedding where progression structure is preserved. This function must be permutation-invariant over neighbors:
\begin{equation}
    f_\theta(x_i, \mathcal{N}_i) = f_\theta(x_i, \pi(\mathcal{N}_i)) \quad \forall \text{ permutations } \pi
\end{equation}
while remaining sensitive to spatial organization. The Set Transformer architecture satisfies this by construction.

\subsection{Data Preprocessing}

\subsubsection{Single-Nucleus RNA Sequencing}

Raw count matrices undergo stringent quality control. We filter cells with fewer than 500 detected genes or greater than 20\% mitochondrial reads—thresholds established in prior lung atlases. Doublets are detected via Scrublet with expected rate calibrated to observed density. Ambient RNA is corrected using SoupX.

Following QC, we normalize using scran pooling-based size factors, robust to snRNA-seq zero-inflation. Counts are log-transformed with pseudocount one:
\begin{equation}
    x_{ig} = \log\left(\frac{c_{ig}}{s_i} + 1\right)
\end{equation}
where $c_{ig}$ is the raw count and $s_i$ the size factor. We select $G = 2000$ highly variable genes via Seurat v3 variance-stabilizing transformation.

Batch effects are corrected using Harmony, which iteratively adjusts embeddings to remove technical variation while preserving biology. Cell type annotations are transferred from HLCA via label propagation on the corrected embedding.

\subsubsection{Spatial Transcriptomics}

Matched Visium data provides spatial context. Since Visium spots (55 $\mu$m diameter) contain multiple cells, we deconvolve spot-level expression into cell type proportions using Tangram with the annotated snRNA-seq as reference.

Spatial coordinates are converted to micrometers using the 100 $\mu$m spot pitch. Neighborhood graphs are constructed by computing pairwise Euclidean distances between spot centroids, retaining edges within 500 $\mu$m maximum radius.

\subsection{Dual-Reference Geometry}

Raw expression vectors live in a space where Euclidean distance poorly reflects biological similarity. We address this through dual-reference mapping onto two complementary atlases: the Human Lung Cell Atlas (HLCA), providing a healthy anchor with 584,944 cells across 58 types, and the Lung Cancer Atlas (LuCA), capturing 790,079 cells representing malignant heterogeneity.

Since reference models expect Ensembl gene identifiers (ENSG) while query data typically uses gene symbols, we first constructed a symbol$\to$ENSG mapping by combining annotations from both references. This reconciliation step ensured maximum gene overlap: HLCA provided mappings for its 2,000 highly variable genes, while LuCA contributed mappings for all 17,764 genes in the Core Atlas. The combined mapping enabled conversion of $>$85\% of query genes to ENSG identifiers.

For each reference, we applied scArches surgical fine-tuning to pretrained scANVI models. The procedure freezes decoder weights $\theta$ while training a query-specific encoder $\phi'$:
\begin{equation}
    \mathcal{L}_{\text{surgery}} = -\mathbb{E}_{q_{\phi'}(z|x)}[\log p_\theta(x|z)] + \text{KL}(q_{\phi'}(z|x) \| p(z))
\end{equation}
Training proceeds for up to 200 epochs with early stopping on validation ELBO (patience 15, 90/10 train-validation split). This yields embeddings $z_i^{\text{HLCA}} \in \mathbb{R}^{30}$ and $z_i^{\text{LuCA}} \in \mathbb{R}^{10}$.

Since these spaces have different scales and semantics, we L2-normalize before fusion. The default concatenation strategy produces:
\begin{equation}
    z_i^{\text{fused}} = \left[ \frac{z_i^{\text{HLCA}}}{\|z_i^{\text{HLCA}}\|_2} \;\Big\|\; \frac{z_i^{\text{LuCA}}}{\|z_i^{\text{LuCA}}\|_2} \right] \in \mathbb{R}^{40}
\end{equation}
The HLCA component captures proximity to healthy states; the LuCA component captures similarity to cancer phenotypes.

\subsubsection{Alternative Fusion Strategies}

Beyond concatenation, we implement two alternative fusion strategies evaluated as ablations:

\paragraph{Learned Weighted Fusion.} A trainable scalar $w \in [0,1]$ interpolates between normalized reference embeddings projected to a common dimension $D_{\text{out}}$:
\begin{equation}
    z_i^{\text{fused}} = w \cdot \text{proj}(\hat{z}_i^{\text{HLCA}}) + (1-w) \cdot \text{proj}(\hat{z}_i^{\text{LuCA}})
\end{equation}
where $\hat{z}$ denotes L2-normalized embeddings and projection handles dimension mismatch via zero-padding. The weight $w$ is initialized to 0.5 and learned during training, allowing the model to discover optimal reference weighting.

\paragraph{Confidence-Weighted Fusion.} Mapping confidence scores dynamically weight each reference per cell:
\begin{equation}
    z_i^{\text{fused}} = \frac{c_i^{\text{HLCA}} \cdot \hat{z}_i^{\text{HLCA}} + c_i^{\text{LuCA}} \cdot \hat{z}_i^{\text{LuCA}}}{c_i^{\text{HLCA}} + c_i^{\text{LuCA}} + \epsilon}
\end{equation}
This allows cells with uncertain mapping to one reference to rely more heavily on the other, potentially improving robustness for transitional phenotypes poorly captured by either atlas alone.

To quantify mapping confidence, we compute calibrated scores via percentile-rank normalization. Let $\delta_i^{\mathcal{R}}$ be the mean distance to 50 nearest reference neighbors:
\begin{equation}
    \delta_i^{\mathcal{R}} = \frac{1}{k} \sum_{j \in \text{kNN}_k(z_i^{\mathcal{R}})} \|z_i^{\mathcal{R}} - z_j^{\mathcal{R}}\|_2
\end{equation}
Raw distances are density-confounded, so they are converted to percentile ranks:
\begin{equation}
    c_i^{\mathcal{R}} = 1 - \frac{\text{rank}(\delta_i^{\mathcal{R}}) - 1}{|\mathcal{R}| - 1}
\end{equation}
This yields confidence in $[0,1]$ comparable across references regardless of local density.

\subsubsection{Post-hoc Confidence Calibration}

Percentile-rank calibration ensures comparability but does not guarantee that confidence scores reflect true accuracy. We apply temperature scaling \cite{Guo2017-calibration} as post-hoc calibration to improve reliability:
\begin{equation}
    \tilde{c}_i = \sigma\left(\frac{\text{logit}(c_i)}{T}\right)
\end{equation}
where $T > 0$ is a learned temperature and $\sigma$ is the sigmoid function. The optimal temperature is found by minimizing Expected Calibration Error (ECE) on a held-out validation set:
\begin{equation}
    \text{ECE} = \sum_{b=1}^{B} \frac{|B_b|}{N} \left| \text{acc}(B_b) - \text{conf}(B_b) \right|
\end{equation}
where samples are binned by confidence into $B=15$ bins, and $\text{acc}(B_b)$, $\text{conf}(B_b)$ are the mean accuracy and confidence within bin $b$. Temperature scaling preserves ranking while improving calibration---a property desirable when confidence scores inform downstream weighting decisions.

\subsection{Receiver-Centered Niche Encoder}

The core contribution is a hierarchical Set Transformer modeling how receiver cells integrate microenvironment signals. Unlike graph neural networks that propagate messages symmetrically, our architecture designates a focal receiver attending to surrounding senders—reflecting the directional nature of ligand-receptor signaling.

\subsubsection{Token Vocabulary}

The model operates over a sequence of $K = R + 4$ tokens encoding receiver identity, spatial structure, and biological context.

\paragraph{Receiver Token.} The receiver token $t_0 \in \mathbb{R}^d$ represents the focal cell. We project the fused embedding through a learned transformation:
\begin{equation}
    t_0 = \text{LayerNorm}(W_{\text{recv}} z_i^{\text{fused}} + b_{\text{recv}})
\end{equation}
where $W_{\text{recv}} \in \mathbb{R}^{d \times 40}$ and $b_{\text{recv}} \in \mathbb{R}^d$. During masked training, $t_0$ is replaced by a learned embedding $\mathbf{m} \in \mathbb{R}^d$.

\paragraph{Spatial Ring Tokens.} The neighborhood is discretized into $R = 4$ concentric rings with boundaries $\mathbf{d} = [0, 50, 100, 200, 500]$ micrometers. These scales correspond to juxtacrine contact ($<50$ $\mu$m), short-range paracrine diffusion (50–100 $\mu$m), medium-range signaling (100–200 $\mu$m), and broader tissue organization (200–500 $\mu$m).

For ring $r$, the neighbor set is $\mathcal{N}_i^{(r)} = \{j : d_{r-1} \leq r_{ij} < d_r\}$. Each ring token aggregates:
\begin{equation}
    t_r = \text{LayerNorm}\left(W_{\text{ring}} \left[ \mathbf{p}_r \;\|\; \mathbf{e}_r \;\|\; \rho_r \;\|\; \mathbf{1}_r \right] + b_{\text{ring}}\right)
\end{equation}
where:
\begin{itemize}
    \item $\mathbf{p}_r \in \mathbb{R}^C$ is the cell type proportion vector with $p_{rc} = |\{j \in \mathcal{N}_i^{(r)} : \tau_j = c\}| / \max(|\mathcal{N}_i^{(r)}|, 1)$
    \item $\mathbf{e}_r \in \mathbb{R}^M$ is mean expression over $M = 8$ gene modules (proliferation, EMT, hypoxia, glycolysis, oxidative phosphorylation, interferon response, TGF-$\beta$ signaling, apoptosis)
    \item $\rho_r = \log(|\mathcal{N}_i^{(r)}| + 1)$ is log-density
    \item $\mathbf{1}_r \in \mathbb{R}^R$ is one-hot ring position
\end{itemize}
The projection $W_{\text{ring}} \in \mathbb{R}^{d \times (C + M + 1 + R)}$ maps to model dimension.

\paragraph{Confidence Token.} Reference mapping confidence provides uncertainty information:
\begin{equation}
    t_{R+1} = \text{LayerNorm}(W_{\text{conf}} [c_i^{\text{HLCA}} \;\|\; c_i^{\text{LuCA}} \;\|\; c_i^{\text{HLCA}} - c_i^{\text{LuCA}}] + b_{\text{conf}})
\end{equation}
The difference term explicitly encodes healthy-vs-cancer reference balance.

\paragraph{Pathway Token.} Pathway activity scores from 50 Hallmark gene sets (computed via AUCell) provide functional context:
\begin{equation}
    t_{R+2} = \text{LayerNorm}(W_{\text{path}} \mathbf{a}_i + b_{\text{path}})
\end{equation}
where $\mathbf{a}_i \in \mathbb{R}^{50}$ contains activity scores.

\paragraph{Stage Token.} When labels are available, a learned stage embedding provides supervision:
\begin{equation}
    t_{R+3} = E_{\text{stage}}[s_i]
\end{equation}
where $E_{\text{stage}} \in \mathbb{R}^{|\mathcal{S}| \times d}$ is a learnable embedding table. During inference or when labels are unavailable, this token is replaced by a learned $[\text{UNK}]$ embedding.

\paragraph{CLS Token.} A learnable classification token $t_{\text{CLS}} \in \mathbb{R}^d$ aggregates sequence-level information, analogous to BERT.

\subsubsection{Neighbor Embedding and Positional Encoding}

Within each ring, individual neighbor cells are embedded for fine-grained attention. For neighbor $j$ in ring $r$:
\begin{equation}
    \tilde{z}_j = W_{\text{in}} z_j^{\text{fused}} + \text{PE}(r_{ij})
\end{equation}
where $W_{\text{in}} \in \mathbb{R}^{d \times 40}$ projects to model dimension.

The positional encoding $\text{PE}(r) \in \mathbb{R}^d$ is sinusoidal over radial distance:
\begin{equation}
    \text{PE}(r)_{2k} = \sin\left(\frac{r}{10000^{2k/d}}\right), \quad \text{PE}(r)_{2k+1} = \cos\left(\frac{r}{10000^{2k/d}}\right)
\end{equation}
This encoding is smooth and unbounded, allowing the model to distinguish fine spatial differences within rings while generalizing to unseen distances.

\subsubsection{Handling Variable-Size Neighborhoods}

Rings contain variable numbers of neighbors—dense tumor cores may have hundreds, sparse stroma only a few. We handle this through sampling, padding, and masking.

When $|\mathcal{N}_i^{(r)}| > n_{\max} = 128$, we uniformly sample $n_{\max}$ neighbors without replacement. When $|\mathcal{N}_i^{(r)}| < n_{\min} = 1$, we pad with a learned $[\text{PAD}]$ embedding $\mathbf{p} \in \mathbb{R}^d$.

Attention masks prevent padded positions from influencing computation:
\begin{equation}
    \text{mask}_{ij} = \begin{cases} 0 & \text{if } j \text{ is valid} \\ -\infty & \text{if } j \text{ is padding} \end{cases}
\end{equation}
This mask is applied additively before softmax, zeroing out attention to padded tokens.

\subsubsection{Induced Set Attention Block (ISAB)}

Processing neighbor sets requires permutation-invariant attention. Standard self-attention has $O(n^2)$ complexity; with neighborhoods of 100+ cells, this becomes prohibitive. We use Induced Set Attention Blocks that mediate interactions through $m$ learned inducing points, achieving $O(nm)$ complexity.

The ISAB consists of two multihead attention operations:
\begin{equation}
    H = \text{MAB}(I, X), \quad \text{ISAB}_m(X) = \text{MAB}(X, H)
\end{equation}
where $I \in \mathbb{R}^{m \times d}$ are learned inducing points (initialized via Xavier) and $X \in \mathbb{R}^{n \times d}$ is the input set.

Each MAB applies pre-norm multihead attention with residual connection:
\begin{equation}
    \text{MAB}(Q, KV) = Q + \text{Dropout}(\text{MHA}(\text{LN}(Q), \text{LN}(KV)))
\end{equation}
followed by a feedforward block:
\begin{equation}
    \text{FFN}(x) = W_2 \cdot \text{GELU}(W_1 \cdot \text{LN}(x) + b_1) + b_2
\end{equation}
with expansion factor 4 (i.e., $W_1 \in \mathbb{R}^{d \times 4d}$, $W_2 \in \mathbb{R}^{4d \times d}$).

The multihead attention splits into $h$ heads:
\begin{equation}
    \text{MHA}(Q, K, V) = \text{Concat}(\text{head}_1, \ldots, \text{head}_h) W^O
\end{equation}
\begin{equation}
    \text{head}_i = \text{softmax}\left(\frac{Q W_i^Q (K W_i^K)^\top}{\sqrt{d/h}} + \text{mask}\right) V W_i^V
\end{equation}
where $W_i^Q, W_i^K, W_i^V \in \mathbb{R}^{d \times d/h}$ and $W^O \in \mathbb{R}^{d \times d}$.

\subsubsection{Pooling by Multihead Attention (PMA)}

To produce fixed-size outputs from variable-size sets, we use PMA with $k$ learned seed vectors $S \in \mathbb{R}^{k \times d}$:
\begin{equation}
    \text{PMA}_k(X) = \text{MAB}(S, X) \in \mathbb{R}^{k \times d}
\end{equation}
The seeds query the set, extracting $k$ summary vectors. For ring aggregation we use $k=1$; for final output we use $k=4$ to capture complementary aspects.

\subsubsection{Hierarchical Architecture}

The full encoder operates in three stages.

\paragraph{Stage 1: Intra-Ring Aggregation.} Each ring's neighbors are processed independently through $L_1 = 2$ ISAB layers:
\begin{equation}
    H_r = \text{ISAB}_m^{(L_1)}(\{\tilde{z}_j : j \in \mathcal{N}_i^{(r)}\})
\end{equation}
Then collapsed via PMA:
\begin{equation}
    \tilde{h}_r = \text{PMA}_1(H_r) \in \mathbb{R}^d
\end{equation}

\paragraph{Stage 2: Cross-Ring Attention.} The receiver token attends to ring summaries and context tokens. We form the memory:
\begin{equation}
    M = [\tilde{h}_1, \ldots, \tilde{h}_R, t_{R+1}, t_{R+2}, t_{R+3}] \in \mathbb{R}^{(R+3) \times d}
\end{equation}
Cross-attention updates the receiver:
\begin{equation}
    \tilde{t}_0 = t_0 + \text{Dropout}(\text{MHA}(\text{LN}(t_0), \text{LN}(M)))
\end{equation}
followed by feedforward:
\begin{equation}
    \hat{t}_0 = \tilde{t}_0 + \text{Dropout}(\text{FFN}(\tilde{t}_0))
\end{equation}

\paragraph{Stage 3: Self-Attention Refinement.} The full token sequence undergoes $L_2 = 2$ standard transformer encoder layers:
\begin{equation}
    \mathcal{T}' = \text{TransformerEncoder}^{(L_2)}([t_{\text{CLS}}, \hat{t}_0, \tilde{h}_1, \ldots, \tilde{h}_R, t_{R+1}, t_{R+2}, t_{R+3}])
\end{equation}
These layers enable higher-order interactions—e.g., how immune infiltration modulates the effect of stromal activation.

\paragraph{Output.} Final PMA with $k=4$ seeds extracts the representation:
\begin{equation}
    C = \text{PMA}_4(\mathcal{T}') \in \mathbb{R}^{4 \times d}
\end{equation}
Flattened and projected:
\begin{equation}
    e_i = W_{\text{out}} \cdot \text{Flatten}(C) + b_{\text{out}} \in \mathbb{R}^D
\end{equation}
where $W_{\text{out}} \in \mathbb{R}^{D \times 4d}$ and $D = 512$.

\subsubsection{Decoder Architecture}

The expression decoder $g_\phi$ maps representations to predicted expression. We use a 3-layer MLP with residual connection:
\begin{equation}
    h_1 = \text{GELU}(W_1 e_i + b_1)
\end{equation}
\begin{equation}
    h_2 = \text{GELU}(W_2 \cdot \text{Dropout}(h_1) + b_2) + W_{\text{skip}} e_i
\end{equation}
\begin{equation}
    \hat{x}_i = W_3 h_2 + b_3
\end{equation}
where $W_1 \in \mathbb{R}^{512 \times D}$, $W_2 \in \mathbb{R}^{512 \times 512}$, $W_3 \in \mathbb{R}^{G \times 512}$, and $W_{\text{skip}} \in \mathbb{R}^{512 \times D}$ is the residual projection.

\subsection{Training Objectives}

We train via self-supervised objectives leveraging spatial structure rather than progression labels.

\subsubsection{Masked Receiver Reconstruction}

The primary objective predicts receiver expression from niche context alone. With the receiver token replaced by mask embedding $\mathbf{m}$, the model must infer cell identity from neighbors.

The loss combines MSE and cosine similarity:
\begin{equation}
    \mathcal{L}_{\text{recon}} = \frac{1}{NG} \sum_{i=1}^{N} \sum_{g=1}^{G} (x_{ig} - \hat{x}_{ig})^2
\end{equation}
\begin{equation}
    \mathcal{L}_{\text{cos}} = 1 - \frac{1}{N} \sum_{i=1}^{N} \frac{x_i^\top \hat{x}_i}{\|x_i\|_2 \|\hat{x}_i\|_2}
\end{equation}
\begin{equation}
    \mathcal{L}_{\text{mask}} = \mathcal{L}_{\text{recon}} + 0.5 \cdot \mathcal{L}_{\text{cos}}
\end{equation}
This directly tests our hypothesis: if niche composition predicts receiver state, progression-relevant signals flow through the microenvironment.

\subsubsection{Contrastive Ranking}

The ranking loss encourages progression-ordered embeddings. We define stage distance:
\begin{equation}
    \Delta(s_i, s_j) = |\text{ord}(s_i) - \text{ord}(s_j)|
\end{equation}
where $\text{ord}: \mathcal{S} \to \{0,1,2,3,4\}$ maps stages to ordinals.

Triplets $(i, j^+, j^-)$ are sampled with $\Delta(s_i, s_{j^+}) \leq 1$ (positive) and $\Delta(s_i, s_{j^-}) \geq 2$ (negative). The margin loss is:
\begin{equation}
    \mathcal{L}_{\text{rank}} = \frac{1}{|\mathcal{P}|} \sum_{(i,j^+,j^-) \in \mathcal{P}} \max(0, \|e_i - e_{j^+}\|_2^2 - \|e_i - e_{j^-}\|_2^2 + \gamma)
\end{equation}
with margin $\gamma = 1.0$. We use semi-hard negative mining for stability.

\subsubsection{Cross-View Consistency}

Augmented views should yield similar representations. Neighbor dropout (rate $p = 0.2$) and coordinate jitter ($\sigma = 10$ $\mu$m) are applied:
\begin{equation}
    \mathcal{L}_{\text{consist}} = \frac{1}{N} \sum_{i=1}^{N} \|e_i - e_i'\|_2^2
\end{equation}
To prevent collapse, we add variance regularization:
\begin{equation}
    \mathcal{L}_{\text{var}} = \frac{1}{D} \sum_{d=1}^{D} \max(0, 1 - \text{Var}(e_{:,d}))
\end{equation}
Combined: $\mathcal{L}_{\text{cv}} = \mathcal{L}_{\text{consist}} + 0.1 \cdot \mathcal{L}_{\text{var}}$.

\subsubsection{Spatial Coherence}

Nearby cells should have similar representations:
\begin{equation}
    \mathcal{L}_{\text{spatial}} = \frac{1}{N} \sum_{i=1}^{N} \sum_{j \in \text{kNN}_{10}(i)} w_{ij} \|e_i - e_j\|_2^2
\end{equation}
with Gaussian weights $w_{ij} = \exp(-r_{ij}^2 / 2\sigma^2)$, $\sigma = 100$ $\mu$m.

\subsubsection{Biological Group Structure}

Cell type annotations provide weak supervision via supervised contrastive loss:
\begin{equation}
    \mathcal{L}_{\text{group}} = -\frac{1}{N} \sum_{i=1}^{N} \frac{1}{|P(i)|} \sum_{j \in P(i)} \log \frac{\exp(\text{sim}(e_i, e_j)/\tau)}{\sum_{k \neq i} \exp(\text{sim}(e_i, e_k)/\tau)}
\end{equation}
where $P(i) = \{j : \tau_j = \tau_i, j \neq i\}$, $\text{sim}$ is cosine similarity, and $\tau = 0.1$.

\subsubsection{Combined Objective}

The total loss combines objectives with configurable weights $\lambda$:
\begin{equation}
    \mathcal{L} = \lambda_{\text{mask}} \mathcal{L}_{\text{mask}} + \lambda_{\text{rank}} \mathcal{L}_{\text{rank}} + \lambda_{\text{cv}} \mathcal{L}_{\text{cv}} + \lambda_{\text{spatial}} \mathcal{L}_{\text{spatial}} + \lambda_{\text{group}} \mathcal{L}_{\text{group}}
\end{equation}

The default configuration weights masked reconstruction at 70\%, emphasizing the primary self-supervised objective:
\begin{equation}
    (\lambda_{\text{mask}}, \lambda_{\text{rank}}, \lambda_{\text{cv}}, \lambda_{\text{spatial}}, \lambda_{\text{group}}) = (0.70, 0.10, 0.10, 0.05, 0.05)
\end{equation}

Alternative weight configurations are evaluated as ablations to assess sensitivity:
\begin{itemize}
    \item \textit{Equal weights}: $(0.20, 0.20, 0.20, 0.20, 0.20)$ --- tests whether the 70\% emphasis on masked reconstruction is necessary
    \item \textit{Masked-only}: $(1.0, 0, 0, 0, 0)$ --- tests whether auxiliary objectives provide meaningful regularization
\end{itemize}
The rationale for emphasizing masked reconstruction is that it directly tests the niche gating hypothesis: if receiver state can be predicted from neighborhood context, progression-relevant information flows through the microenvironment.

\subsection{Baseline Architectures}

We compare against a ladder of increasing structural sophistication.

\textbf{Pooling MLP.} Concatenates receiver embedding with mean/max statistics:
\begin{equation}
    e_i = \text{MLP}([z_i^{\text{fused}} \| \text{mean}(\mathcal{N}_i) \| \text{max}(\mathcal{N}_i)])
\end{equation}
No set structure or permutation invariance.

\textbf{DeepSets.} Permutation-invariant via sum aggregation:
\begin{equation}
    e_i = \rho\left(z_i^{\text{fused}}, \sum_{j \in \mathcal{N}_i} \phi(z_j^{\text{fused}})\right)
\end{equation}
Cannot model neighbor interactions.

\textbf{Flat Set Transformer.} ISAB over full neighbor set without ring hierarchy:
\begin{equation}
    e_i = \text{PMA}_k(\text{ISAB}^{(L)}(\{z_i^{\text{fused}}\} \cup \mathcal{N}_i))
\end{equation}
Tests whether explicit spatial scales help.

\textbf{GraphSAGE.} Message passing with learned aggregation:
\begin{equation}
    h_i^{(l+1)} = \sigma\left(W^{(l)} [h_i^{(l)} \| \text{AGG}(\{h_j^{(l)} : j \in \mathcal{N}_i\})]\right)
\end{equation}
Symmetric propagation unlike our receiver-centered design.

\textbf{GAT.} Attention-weighted message passing:
\begin{equation}
    h_i^{(l+1)} = \sigma\left(\sum_{j \in \mathcal{N}_i} \alpha_{ij}^{(l)} W^{(l)} h_j^{(l)}\right)
\end{equation}
Still symmetric—all nodes are both senders and receivers.

\subsection{Ablation Studies}

Ablations are organized in three tiers by centrality to claims.

\textbf{Tier 1: Core Hypothesis.}
\begin{itemize}
    \item \textit{No-niche}: Receiver embedding only, $e_i = \text{MLP}(z_i^{\text{fused}})$
    \item \textit{Pooled-niche}: Mean pooling instead of attention
\end{itemize}

\textbf{Tier 2: Architecture.}
\begin{itemize}
    \item \textit{Flat attention}: No ring hierarchy
    \item \textit{No inducing points}: Standard $O(n^2)$ attention
    \item \textit{Single-head}: $h=1$ instead of $h=8$
    \item \textit{No positional encoding}: Remove sinusoidal PE
\end{itemize}

\textbf{Tier 3: Reference Geometry.}
\begin{itemize}
    \item \textit{HLCA-only}: Healthy reference only
    \item \textit{LuCA-only}: Cancer reference only
    \item \textit{Unnormalized}: Concatenate without L2 normalization
    \item \textit{Learned fusion}: Trainable weight instead of concatenation
    \item \textit{Weighted fusion}: Confidence-weighted combination
\end{itemize}

\textbf{Tier 4: Training Objectives.}
\begin{itemize}
    \item \textit{Equal loss weights}: All objectives weighted equally at 0.20
    \item \textit{No auxiliary losses}: Only masked reconstruction (weight 1.0)
\end{itemize}

\textbf{Tier 5: Confidence Calibration.}
\begin{itemize}
    \item \textit{No calibration}: Raw percentile-rank confidence without temperature scaling
    \item \textit{Temperature calibration}: Post-hoc temperature scaling to minimize ECE
\end{itemize}

\subsection{Evaluation Metrics}

\textbf{Reconstruction.} MSE and Pearson correlation between predicted and true expression.

\textbf{Progression Discrimination.} Linear probe on frozen embeddings: macro F1 and ordinal accuracy (correct if within 1 stage).

\textbf{Biological Coherence.} Silhouette score and adjusted Rand index over cell type annotations.

\textbf{Generalization.} All metrics computed on held-out donors; train-test gap reveals overfitting.

\subsection{Implementation Details}

\subsubsection{Architecture Hyperparameters}

Model dimension $d = 256$. Inducing points $m = 32$. Attention heads $h = 8$. Intra-ring ISAB layers $L_1 = 2$. Cross-ring transformer layers $L_2 = 2$. PMA seeds $k = 4$. Output dimension $D = 512$. Dropout rate 0.1 throughout.

\subsubsection{Initialization}

Linear projections are initialized with Xavier and embeddings (inducing points, seeds, special tokens) with $\mathcal{N}(0, 0.02)$. Layer norm parameters are set to $\gamma = 1$, $\beta = 0$.

\subsubsection{Optimization}

AdamW optimizer with $\beta_1 = 0.9$, $\beta_2 = 0.999$, weight decay 0.01. Base learning rate $10^{-4}$ with linear warmup over 1000 steps, then cosine decay to $10^{-6}$:
\begin{equation}
    \eta_t = \eta_{\min} + \frac{\eta - \eta_{\min}}{2}\left(1 + \cos\left(\frac{t - t_{\text{warmup}}}{T - t_{\text{warmup}}}\pi\right)\right)
\end{equation}

Batch size 256 neighborhoods. Gradient accumulation over 4 steps when memory-constrained (effective batch 1024). Gradient clipping at norm 1.0.

\subsubsection{Early Stopping}

Training proceeds for up to 100 epochs. Validation loss (combined objective on held-out donors) is monitored; training halts after 10 epochs without improvement. The best checkpoint by validation loss is retained.

\subsubsection{Data Splitting}

Donor-held-out 5-fold cross-validation is employed, ensuring all cells from a donor appear exclusively in train or validation. Folds are stratified by stage distribution, and no spatial coordinate overlap between splits is permitted.

\subsubsection{Inference}

At inference, the receiver token is \textit{not} masked—we use the actual receiver embedding. The model outputs the representation $e_i$ for downstream tasks. Attention weights from cross-ring attention provide interpretability, quantifying which spatial scales and niche components influenced the representation.

\subsubsection{Computational Resources}

NVIDIA H100 80GB GPUs with mixed-precision (FP16) training. Reference mapping: ~4 hours for 787K cells. Encoder training: ~8 hours for 100 epochs. Full ablation suite: ~80 GPU-hours.
